\documentclass[conference]{IEEEtran}
\usepackage{cite,amsmath,amssymb,graphicx,booktabs}
\usepackage[hidelinks]{hyperref}
\begin{document}
\title{READER: Observed-Prefix Reverse Reading for Causal EEG Decoding}
\input{authors.tex}
\maketitle
\begin{abstract}
EEG decisions made before trial completion require representations that respect
the available observation boundary. We study READER, which combines a causal
temporal--spatial front end with forward and reverse temporal convolutional
readers. Building on an existing forward--reverse readout and fusion design,
READER recomputes its reverse summary from the observed token prefix. A static
feature-wise gate combines the latest-token forward state with a start-anchored
reverse summary. Training uses only the full-trial endpoint objective, and the
same checkpoint is evaluated at shorter observation lengths.
\input{abstract_results.tex}
\end{abstract}
\begin{IEEEkeywords}
EEG classification, causal decoding, prefix prediction, temporal convolutional networks, brain--computer interface
\end{IEEEkeywords}

\section{Introduction}
Electroencephalography (EEG) provides noninvasive measurements of brain activity
with high temporal resolution. Deep EEG classifiers combine temporal and spatial
filtering, attention, and time--frequency representations
\cite{lawhern2018eegnet,altaheri2022physics,song2022eeg,hong2025bidirectional}.
In brain--computer interfaces, full-trial classification and decisions during
acquisition are different evaluation settings: a strong endpoint classifier
need not perform well when less EEG is available.

Causal temporal convolutions are useful EEG modeling components
\cite{salami2022eeg}. Causality of one block, however, does not establish that
preceding feature extraction and the final readout obey the same observation
boundary. Conversely, reverse-order processing of an already observed segment
is compatible with a decision at that segment's endpoint. The constraint is
information availability, not processing direction. A reverse representation
computed from the complete trial cannot generally be reused unchanged for an
earlier decision.

Building on BiTE's forward--reverse readout and its released implementation's
feature-wise fusion \cite{hong2025bidirectional,hong2025bitecode}, READER studies
a causal tokenizer, repeated reverse encoding of the observed prefix, and
endpoint-only training. Both readers receive the same observations. The
hypothesized benefit concerns their different nonlinear processing orders and
readout anchors, not additional EEG information or a claimed physiological
mechanism. A forward-only model also supports prefix predictions; therefore,
checkpoint reuse alone is not evidence for the value of reverse reading.

We evaluate this question in a fresh matched within-subject study on BCICIV-2A,
BCICIV-2B, and SD-SSVEP. Four architectures share the same front end and training
recipe: Compact, READER, Forward--Forward (FF), and Compact--Mean. BiTE is
retrained as a separate endpoint comparator. Existing numerical tables from
earlier model variants are not combined with the new cohort.

The contributions are: (1) a specified observed-prefix EEG formulation with an
explicit prepared-input inference boundary; (2) endpoint-trained comparisons at
identical observation lengths, including capacity and readout controls; and
(3) finite trained-model causality tests, paired subject-level uncertainty,
and scoped inference-cost measurements.

\section{Related Work}
\subsection{Temporal and Bidirectional EEG Decoding}
EEGNet combines temporal, depthwise spatial, and separable convolutional
filtering \cite{lawhern2018eegnet}. EEG-ITNet uses inception modules and dilated
temporal convolutions \cite{salami2022eeg}; ATCNet combines convolution,
attention, and temporal modeling \cite{altaheri2022physics}. EEG Conformer
combines local convolutional features with self-attention
\cite{song2022eeg}. These architectures motivate rich EEG representations,
but full-segment performance does not characterize their earlier predictions.

BiTE combines temporal and STFT streams, pyramid time--frequency attention,
and bidirectional temporal convolution \cite{hong2025bidirectional}. Its final
reverse readout already corresponds to the original first position, and its
released implementation uses feature-wise sigmoid fusion
\cite{hong2025bitecode}. READER does not claim these readout anchors or this
fusion rule as new. It studies their combination with prefix-causal tokenization
and repeated inference restricted to the available EEG.

Prefix recomputation is also used outside EEG, for example in re-translation for
simultaneous machine translation \cite{sen2021university}. The present study
addresses the usefulness of a particular EEG representation rather than
claiming recomputation as a new general principle.

\subsection{Early and Anytime EEG Decisions}
Short-window EEG classifiers address decisions from limited observations
\cite{hong2022deep}. Dynamic stopping instead decides whether to issue an output
or acquire more evidence \cite{nakanishi2015dynamic,ahmadi2024bayesian}.
Here, \emph{anytime} denotes evaluating one trained model at several observation
boundaries. It does not imply calibrated confidence, an optimized stopping
rule, or a demonstrated acquisition-to-action latency advantage.

An endpoint-trained forward-only checkpoint is the direct comparison for
READER's reuse setting. A deadline-specific specialist would address a different
training condition. The current fresh study focuses on matched reused
checkpoints and does not reuse previously reported specialist or LOSO scores.

\section{Method}
\label{sec:method}
\subsection{Architecture and Observation Boundary}
\begin{figure*}[t]
\centering
\includegraphics[width=.94\textwidth]{Fig.png}
\caption{READER at an observed token prefix. A causal temporal--spatial front
end produces positioned tokens. The forward reader returns its latest-token
state; the reverse reader processes only the supplied prefix in reverse and
returns the state corresponding to original token 1. A static feature-wise
gate and shared classifier combine the two representations. The classifier is
supervised only at the full-trial endpoint.}
\label{fig:reader}
\end{figure*}
Let $\mathbf X\in\mathbb R^{C\times T}$ be a prepared EEG input and
$m\in\{1,\ldots,T\}$ the available sample count. The Reverse-Enabled Anytime
Decoder for EEG Recognition (READER) contains a causal front end and two
independently parameterized readers (Fig.~\ref{fig:reader}). Model-level
causality is defined on the prepared-input timeline; it does not establish
causality of acquisition or the earlier raw-data preparation.

\subsection{Causal Temporal--Spatial Front End}
Three stride-one temporal convolutions have kernel lengths 16, 32, and 64,
16 output maps each, no biases, and zero left padding of $k-1$ samples.
Concatenation yields 48 temporal carrier maps followed by BatchNorm.
A bias-free depthwise spatial convolution across all $C$ channels learns two
filters per carrier, yielding 96 features. BatchNorm and LeakyReLU with slope
0.2 follow. Each spatial filter's weights are constrained to $\ell_2$ norm 1.

Let $\mathbf u_q\in\mathbb R^{96}$ denote the feature at sample $q$.
Non-overlapping mean pooling uses width $p=32$ for 2A and 2B and $p=4$
for SD-SSVEP. Define $N_m=\lceil m/p\rceil$ and $s_j(m)=\min(jp,m)$.
The pooled token input is
\begin{equation}
\bar{\mathbf u}_j^{(m)}=
\frac{\sum_{q=(j-1)p+1}^{s_j(m)}\mathbf u_q}{s_j(m)-(j-1)p},
\quad j=1,\ldots,N_m.
\end{equation}
Thus, incomplete windows average only observed samples. For $m=1000$, $p=32$,
the final window contains eight samples. At 250, 500, and 750 samples it
contains 26, 20, and 14 samples, respectively. A completed token computed from
a longer input cannot substitute for one of these partial tokens.

Dropout 0.3 precedes a biased linear projection from 96 to 64 dimensions.
Learned positional vectors, initialized from $\mathcal N(0,0.02^2)$, are added
at the original trial positions. The table has 32 positions for 2A/2B and 64
for SD-SSVEP. Write the resulting sequence as
$\mathbf z_1^{(m)},\ldots,\mathbf z_{N_m}^{(m)}$. For a fixed supplied prefix,
we omit the superscript and set $t=N_m$. Distinct sample counts can yield
the same token count but different final-token contents.

\subsection{Forward and Reverse Readers}
Each reader has three residual TCN blocks with dilations 1, 2, and 4.
A block has two bias-free depthwise kernel-6 convolutions with zero left
padding of $5d$ samples. Each convolution is followed by affine BatchNorm,
LeakyReLU(0.2), and dropout 0.3. The block adds its residual output to its
input with no subsequent activation. First-convolution weights use Xavier
uniform initialization; second-convolution weights and BatchNorm offsets are
zero initially. BatchNorm scales start at one. Blocks therefore begin as
identity mappings. All BatchNorm layers use $\epsilon=10^{-5}$ and momentum 0.1.

With stride one and fixed evaluation statistics, each reader's theoretical
receptive field is $R=1+2(6-1)(1+2+4)=71$ tokens. This covers the evaluated
32- and 64-token sequences; it indicates architectural access, not learned
use of every token. The forward and reverse summaries are
\begin{align}
\mathbf f_t&=[\mathcal F(\mathbf z_1,\ldots,\mathbf z_t)]_t,\\
\mathbf b_t&=[\mathcal B(\mathbf z_t,\ldots,\mathbf z_1)]_{\mathrm{last}}.
\end{align}
Reversal follows positional encoding, preserving original positional vectors.
The reverse readout is anchored at original token 1 and initially equals
$\mathbf z_1$. For longer sequences, it is limited to its first-$R$-token
original-order support. The evaluated formulation therefore concerns bounded
trial prefixes rather than unrestricted continuously growing histories.

\subsection{Fusion and Controls}
Following the feature-wise fusion in the released BiTE implementation
\cite{hong2025bitecode}, a static gate combines the states:
\begin{equation}
\mathbf g=\sigma(\mathbf a),\qquad
\mathbf h_t=(1-\mathbf g)\odot\mathbf f_t+\mathbf g\odot\mathbf b_t.
\end{equation}
The 64 gate parameters start at zero; gates start at 0.5 and are shared across
trials and observation lengths. A biased linear classifier maps $\mathbf h_t$
to $K$ class logits. Its weight rows have maximum $\ell_2$ norm 0.25.
Spatial and classifier constraints are projected after optimizer updates.
Inference for the four controlled models does not mutate these weights.

Compact uses $\mathbf f_t$ without a second reader or gate. FF uses a second
independent forward TCN and the same fusion rule. Compact--Mean instead uses
$t^{-1}\sum_{j=1}^{t}[\mathcal F(\mathbf z_1,\ldots,\mathbf z_t)]_j$.
Its equal-token mean includes the partial final token as one token, and the
whole model is trained with this readout. Shared parameters initialize
identically across all four models for each seed. The additional READER and FF
reader plus gate contains 3,136 parameters. FF controls capacity and fusion,
but changes direction and readout anchor together, so it does not perfectly
isolate temporal order from all readout effects.

\subsection{Inference Causality and Cost}
Let $\Phi_m$ tokenize the supplied $m$ samples and $\mathcal D_\theta$
include the readers, fusion, and classifier. With dropout disabled and
BatchNorm statistics fixed,
\begin{equation}
\boldsymbol\ell_m(\mathbf X)=\mathcal D_\theta(\Phi_m(\mathbf X_{:,1:m})).
\end{equation}
Two prepared inputs equal through $m$ yield the same available sample
features, pooled tokens, positioned sequences, and logits. Completed tokens
may also be reused from a longer causal front-end pass. At a partial boundary,
pooling is recomputed from the available features. Training-time BatchNorm is
not claimed to implement online learning. Raw preprocessing causality and
acquisition-to-output delay are outside this argument.

One reverse pass is linear in the available token count for fixed model width
and depth. Recomputing every growing prefix requires $O(N^2)$ total
reverse-reader work in this implementation. The independent curve routine
shares sample features and completed forward states where possible, but does
not claim constant-cost reverse updates.

\section{Experimental Design}
\subsection{Data and Frozen Preparation}
The study includes BCICIV-2A (9 subjects, 22 channels, 4 classes), BCICIV-2B
(9 subjects, 3 channels, 2 classes), and SD-SSVEP (10 subjects, 8 channels,
12 classes) \cite{tangermann2012review,nakanishi2015comparison}.
Prepared 2A and 2B inputs contain 1000 samples at 250 Hz; SSVEP inputs contain
256 samples at 256 Hz. The original 2A training/evaluation sessions are kept
separate, and 2B sessions 1--3 and 4--5 form the respective roles.
SD-SSVEP retains the supplied per-subject training/test artifacts without
creating a new split. File names, array shapes, retained row ordinals, and
checksums are recorded. These ordinal identifiers do not independently prove
original trial/block independence.

The loader reads the existing prepared NPZ arrays, retains the first 1000 MI
samples when an inclusive endpoint is present, and applies no new band-pass
filter, rereferencing, or covariance alignment. A separate StandardScaler is
fitted across training trials at every channel--sample coordinate and then
applied unchanged to testing. Its means and scales are persisted, and shorter
prefixes use only the corresponding coordinates. No test-trial statistics
are fitted. Exact prepared-array duplicates across training/test roles are
checked; that test does not establish all properties of the earlier preparation.
All observation durations are relative to the prepared crop, not stimulus
onset or system latency. The inherited raw-to-prepared provenance requires
independent documentation; this study's causal claim starts at the prepared
input.

\subsection{Fresh Matched Training}
This is a new manuscript-based implementation and refit, not a numerical
reproduction of previously recovered checkpoints. Prior benchmark outcomes
informed development; the follow-up does not make those test sets newly
unseen. The new recipe and observation grid are fixed before the follow-up
outputs, and no new test-epoch selection is used.

For each subject, seeds 2025, 2026, and 2027 train Compact, READER, FF,
Compact--Mean, and BiTE: $28\times3\times5=420$ training runs.
All use 600 epochs, Adam (learning rate 0.002, weight decay 0.002,
$\beta=(0.9,0.999)$, $\epsilon=10^{-8}$), batch size 64, label smoothing 0.1,
global gradient-norm clipping at 5, and cosine decay to zero. The final short
batch is retained; there is no data augmentation or intermediate loss.
The scheduler advances after each epoch. FP32 is used with TF32 disabled,
deterministic algorithms requested, and environment metadata recorded.

A separate CPU generator determines the minibatch permutation at epoch $e$
from seed $k+20000+e$ and its hash is compared across models. Training RNGs
are reset after construction using $k+10000$. The extra READER/FF branch is
constructed in an isolated CPU RNG context with fixed seed 918273; it does
not vary independently across the three seeds. Shared initialization hashes
are checked. Dropout masks are not claimed to be identical across models.
Each subject/seed group saves all five final checkpoints before the new
held-role evaluation. Reloading checks use training examples.

BiTE uses the authors' released class at commit
\texttt{924eb32241ba1a7c80dbc4ba097f8c979da17578}
\cite{hong2025bitecode}. Its architecture-specific defaults are retained,
with temporal/STFT kernel and pooling settings 64/32 for MI and 32/8 for
SSVEP. Spectral features use magnitude STFT with a Hann window, hop 1,
centered constant padding, and retained bins specified by the implementation's
4--40-Hz or 8--64-Hz configuration. BiTE is the full-trial endpoint comparator;
no unsupported shorter-input behavior is assumed. It is retrained under the
same optimization recipe, not quoted from the publication.

\subsection{Exact-Duration Evaluation and Uncertainty}
The four controlled models reuse their own endpoint-trained checkpoints.
For MI, the primary grid is 256 through 992 samples in steps of 32, augmented
with 250, 500, 750, and 1000: 28 points covering 1--4 s. SSVEP uses 64 through
256 samples in steps of 4, covering 0.25--1 s. Inputs are truncated before
feature extraction; no longer-input final token is substituted.

Accuracy is computed per seed, seeds are averaged within each subject, and
subjects are averaged within each dataset. Seeds are not prediction ensembles
or independent participants. We report subject-level SDs and paired
95\% percentile intervals from 10,000 subject bootstrap samples. Each
resample retains the corresponding models and all observation lengths.
Intervals are conditional on the fitted runs and are pointwise, not
simultaneous bands or multiplicity-adjusted superiority tests.

For each subject, normalized accuracy--duration area is
\begin{equation}
Q=\frac{1}{\tau_J-\tau_1}\sum_{j=1}^{J-1}
(\tau_{j+1}-\tau_j)\frac{\bar A(\tau_j)+\bar A(\tau_{j+1})}{2}.
\end{equation}
It retains accuracy units and is not ROC AUC. The primary contrast is
READER minus Compact on BCICIV-2A over 1--4 s; other datasets and mechanism
contrasts contextualize the result. No favorable time point is selected
post hoc as the primary outcome.

\subsection{Causality, Diagnostics, and Timing}
Every trained controlled-model checkpoint is tested on four label-independent
prepared examples and four synthetic inputs. The independent full-input
curve routine is compared with direct truncated inference. Future samples
are perturbed, repeated and batch-context inference are checked, and second-reader
input lengths are traced. Partial boundaries are included. Tests are repeated
on in-memory copies with nonzero randomized temporal paths and reset
BatchNorm affine/statistical values. Original checkpoints are not modified.
Maximum absolute logit error must not exceed $10^{-4}$, and a deliberately
future-dependent positive control must exceed it. The endpoint has no future
suffix and is not counted as a future-perturbation success.

Supplementary READER diagnostics record gates, state norms, branch similarity,
and in-prefix interventions. Forward-only and reverse-only predictions using
the jointly trained head are labelled post-hoc interventions, not trained
baselines. Gate magnitude alone is not interpreted as feature importance.

Timing uses subject 1/seed 2025 per dataset, batch size one, ten warm-ups,
50 repetitions, and accelerator synchronization. Median and 95th-percentile
milliseconds are recorded for one requested prefix and both naïve and
shared-feature full-curve routines. Inputs are already normalized and on the
device; BiTE's STFT is included. Disk I/O, transfers, earlier preprocessing,
and acquisition delay are excluded. Timings are tied to the recorded device,
not pooled indiscriminately across different GPU types.

\section{Results}
\subsection{Endpoint Classification}
\begin{table*}[t]
\caption{Fresh within-subject endpoint accuracy (\%). Mean $\pm$ SD is across
seed-averaged subject values; no historical results are mixed in.}
\centering\small\input{endpoint_table.tex}
\label{tab:endpoint}
\end{table*}
\input{endpoint_results.tex}

\subsection{Matched Prefix and Mechanism Comparisons}
\begin{table}[t]
\caption{BCICIV-2A exact-duration accuracy (\%). Each model reuses its own
endpoint-trained checkpoint. FF is Forward--Forward; C--Mean is Compact--Mean.}
\centering\scriptsize\setlength{\tabcolsep}{3pt}\input{2a_prefix_table.tex}
\end{table}
\input{prefix_results.tex}
\begin{figure}[t]
\centering\includegraphics[width=\columnwidth]{2a_reader_minus_compact.png}
\caption{Paired BCICIV-2A READER--Compact accuracy differences on the fixed
common grid. Shading gives pointwise 95\% subject-bootstrap intervals,
not simultaneous confidence bands.}
\end{figure}
\begin{table*}[t]
\caption{READER minus each comparator. Paired mean differences and 95\%
subject-bootstrap intervals. $Q$ is normalized accuracy--duration area;
dataset-specific intervals are not pooled.}
\centering\small\input{mechanism_table.tex}
\label{tab:mechanism}
\end{table*}
The controlled contrasts in Table~\ref{tab:mechanism} evaluate alternative
capacity and readout explanations. A difference from FF does not by itself
isolate reversal from its changed anchor. Per-trial predictions are retained in the experiment archive.
2B and SSVEP curves and all operating points accompany the results.

\subsection{Implementation and Computational Checks}
\begin{table*}[t]
\caption{Finite causality-check maxima across the fresh cohort.
Nondegenerate refers to an in-memory randomized temporal-path copy.
These are implementation checks, not proof of raw preprocessing causality.}
\centering\small\input{causality_table.tex}
\end{table*}
All required checks for the generated cohort satisfied the stated numerical
criterion; trace and positive-control outcomes are recorded in the per-run
reports.
\begin{table}[t]
\caption{Parameter counts and measured batch-one inference time for subject 1,
seed 2025 of each dataset. Times are per-dataset device-specific medians, not
acquisition latency. Curve times use the shared-feature routine; BiTE is
measured only at its full-trial endpoint. Devices, 95th percentiles and naive
recomputation times are retained in the run records.}
\centering\scriptsize\setlength{\tabcolsep}{2pt}\input{timing_table.tex}
\end{table}
Parameter counts and device-specific timing details are provided with the
reproducibility outputs. The formulation does not claim a constant-cost
streaming or an acquisition-to-output latency advantage.

\section{Discussion}
The study examines whether two summaries of the same available EEG improve
prefix decoding under endpoint-only supervision. Endpoint, prefix, and
mechanism comparisons answer different questions. Numerical endpoint gains
alone do not establish improved intermediate decisions. Likewise, availability
of multiple outputs from one checkpoint is not unique to reverse processing.
The paired curve-level result and trained controls determine how strongly
the usefulness of the complete configuration is supported.

\input{mechanism_interpretation.tex}

The FF comparison controls added reader capacity and fusion but changes the
second readout anchor together with order. Compact--Mean tests an alternative
forward summary but is not an exhaustive search over forward-only decoders.
Consequently, the study does not identify a direction-only physiological
mechanism. Shorter inputs also introduce length and boundary-condition shifts
relative to endpoint training; their causal role in accuracy changes is not
separately established by this study.

Subject numbers are small and bootstrap intervals are conditional on the
observed fitted models. The historical benchmark exposure is disclosed rather
than treated as a new untouched validation set. Original raw-to-prepared
provenance and trial/block independence are not inferred from array shapes
or absence of exact duplicates. Causality checks concern the prepared model
input; stopping policy, calibration, unrestricted history length, cross-subject
transfer, and full online deployment are outside the primary evaluation.

\section{Conclusion}
READER combines a causal EEG tokenizer with forward and observed-prefix reverse
readers and reuses an endpoint-trained classifier across observation lengths.
\input{conclusion_results.tex}
The fresh study compares identical input durations, additional-capacity and
readout alternatives, subject-level uncertainty, finite causality checks,
and measured inference costs. Its empirical conclusions concern the complete
bounded-prefix configuration and must be read with the paired contrasts and
scope limitations, not as a universal directionality or efficiency advantage.

\bibliographystyle{IEEEtran}
\bibliography{references}
\end{document}
